\input{../tex-inputs/latex-leseansicht-vorspann}

\section[Arthur Schnitzler an Gustav Schwarzkopf, 20. 6. 1928]{L04478 Arthur Schnitzler an Gustav Schwarzkopf, 20. 6. 1928}
\nopagebreak\mylabel{L04478v}
\rehead{ }\normalsize\beginnumbering\briefempfaengerindex{Schwarzkopf, Gustav@\textsc{Schwarzkopf, Gustav}!zzzSchnitzler, Arthur@\emph{von Arthur Schnitzler}!1928-06-201@{20. 6. 1928}|(be}
\toendnotes[C]{\smallbreak\pagebreak[2]}
\correspDesc{Versand  durch Arthur Schnitzler am 20. 6. 1928 in Bad Ischl
\newline{}Übermittlung  am 20. 6. 1928 in Bad Ischl
\newline{}Erhalt  durch Gustav Schwarzkopf im Zeitraum [21. 6. 1928 – 25. 6. 1928?] in Wien}\toendnotes[C]{\smallbreak}
\Standort{DLA, A:Schnitzler, HS.1985.1.1897.}
\physDesc{Bildpostkarte, 487 Zeichen
\newline{}Handschrift: Bleistift, lateinische Kurrent
\newline{}Versand: 1) Stempel: »\nobreak{}\oindex{Bad Ischl@\textbf{Bad Ischl}|pwk}Bad Ischl, 2\textcolor{gray}{0}. VI. 28, 1\textcolor{gray}{×}\nobreak{}«.   2) Stempel: »\nobreak{}Gebrauchet die heilkräftigen Solbäder Bad Ischl’s\oindex{Bad Ischl@\textbf{Bad Ischl}|pw}\nobreak{}«. }\toendnotes[C]{\smallbreak}\pstart{}{\pb}Hr Guſtav Schwarzkopf\pend{}\pstart{}Wien I\oindex{I., Innere Stadt@\textbf{I., Innere Stadt}|pw}\pend{}\pstart{}Tiefer Graben 17\oindex{Wien@\textbf{Wien}!I., Innere Stadt@\textbf{I., Innere Stadt}!Tiefer Graben 17@\textbf{Tiefer Graben 17}|pw}\pend{}{\bigskip}
\pstart
           \noindent{}\centering{}{\pb}\textcolor{gray}{\textbf{41. Bad Ischl\oindex{Bad Ischl@\textbf{Bad Ischl}|pw} g d. Dachstein\oindex{Dachstein@\textbf{Dachstein}|pw}.}}\pend
           \vspace{1em}
\pstart
           \raggedleft{}{\pb}Ischl\oindex{Bad Ischl@\textbf{Bad Ischl}|pw}{ }20/6 928\pend
           \vspace{0.5em}
\pstart
           lieber Guſtav, so leid es mir thut
               dß Sie nicht geko{\geminationm}en sind –
      ich ka{\geminationn} es nur für mich – nicht für
      Sie bedauern – de{\geminationn} seit 8 Tagen
      regnet es fast ununterbrochen –
      wie es wirklich nur hier regnen kann.
      Trotzdem ist der Aufenthalt angenehm
      und das gute Hotel\oindex{Hotel Kaiserin Elisabeth [Bad Ischl]@\textbf{Hotel Kaiserin Elisabeth [Bad Ischl]}|pwv}, der unverwüstliche
               Zauner\oindex{Konditorei Zauner@\textbf{Konditorei Zauner}|pw} und liebenswürdge Gesellschaft
            {\pb}lassen Verdrossenheit nicht aufko{\geminationm}en.\pend
           
\pstart
           Freitag hoff ich wieder zu Hause zu sein.\pend
           
\pstart
           Ich grüße Sie herzlich, Fr. P.\pwindex{Pollaczek, Clara Katharina 15.\,1.\,1875 Wien – 22.\,7.\,1951 ebd.@\textsc{Pollaczek, Clara Katharina} (15.\,1.\,1875 Wien – 22.\,7.\,1951 ebd.)|pw} desgleichen{\\[\baselineskip]}Ihr \spacefill\mbox{A.}\pend
           \leftskip=0em{}\selectlanguage{ngerman}\endnumbering\briefempfaengerindex{Schwarzkopf, Gustav@\textsc{Schwarzkopf, Gustav}!zzzSchnitzler, Arthur@\emph{von Arthur Schnitzler}!1928-06-201@{20. 6. 1928}|)be}\mylabel{L04478h}
\begin{anhang}
\end{anhang}\newcommand{\dateiname}{L04478}\newcommand{\titel}{Arthur Schnitzler an Gustav Schwarzkopf, 20. 6. 1928}\newcommand{\editorInnen}{Herausgegeben von Jahnke, SelmaMüller, Martin Anton}\input{../tex-inputs/latex-leseansicht-abspann}
